\section{Periodic notation, matrix convention과 단위}
\label{sec:notation}

\subsection{Lattice와 reciprocal lattice}

Primitive lattice matrix를 column convention으로
\begin{equation}
  \mathbf A=(\mathbf a_1\;\mathbf a_2\;\mathbf a_3)
\end{equation}
라 하고 cell volume을
\begin{equation}
  \Omega=\det\mathbf A>0
  \label{eq:cell-volume}
\end{equation}
로 둔다. Reciprocal lattice matrix는
\begin{equation}
  \mathbf B=2\pi\mathbf A^{-\Herm},
  \qquad
  \mathbf a_i\cdot\mathbf b_j=2\pi\delta_{ij}
  \label{eq:reciprocal-lattice}
\end{equation}
이다. Lattice translation은 $\vT=\mathbf A\mathbf n$,
$\mathbf n\in\mathbb Z^3$이고 reciprocal vector는 $\vG=\mathbf B\mathbf m$다.

Cell 안의 적분은 $\int_\Omega\dd\mathbf r$로 표기한다. Fourier pair는 implementation의
FFT normalization에 따라 상수 위치가 달라질 수 있으므로, 이 문서에서는
\begin{equation}
  \rho(\vG)=\int_\Omega e^{-i\vG\cdot\mathbf r}\rho(\mathbf r)\dd\mathbf r,
  \qquad
  \rho(\mathbf r)=\frac{1}{\Omega}\sum_{\vG}
  e^{i\vG\cdot\mathbf r}\rho(\vG)
  \label{eq:fourier-pair}
\end{equation}
를 사용한다.

\subsection{Bloch AO와 k-point weight}

원자 중심 Gaussian $\chi_\mu$로부터 Bloch AO를
\begin{equation}
  \phi_{\mu\vk}(\mathbf r)
  =\frac{1}{\sqrt{N_{\mathrm c}}}
   \sum_{\vT}e^{i\vk\cdot\vT}
   \chi_\mu(\mathbf r-\vT)
  \label{eq:bloch-ao}
\end{equation}
로 정의한다. $N_{\mathrm c}$는 Born--von Karman supercell의 cell 수이며 최종 intensive
quantity에서는 소거된다. Bloch orbital은
\begin{equation}
  \psi_{p\vk}(\mathbf r)
  =\sum_\mu C_{\mu p}^{\vk}\phi_{\mu\vk}(\mathbf r)
  \label{eq:bloch-mo}
\end{equation}
이다.

일반 weight $w_\vk$는 $\sum_\vk w_\vk=1$로 normalization한다. 이 package의 regular
mesh는 equal weight를 사용하므로 $w_\vk=1/N_k$다. 아래 식에서 $1/N_k$가 보이면
equal-weight convention을 뜻한다. Gamma-only도 별도 scalar theory가 아니라
$N_k=1$, $\vk=\mathbf0$인 같은 tensor convention으로 다룬다.

AO overlap과 one-electron matrix는
\begin{align}
  S_{\mu\nu}^{\vk}
  &=\langle\phi_{\mu\vk}|\phi_{\nu\vk}\rangle,\\
  h_{\mu\nu}^{\vk}
  &=\langle\phi_{\mu\vk}|\hat h|\phi_{\nu\vk}\rangle
  \label{eq:ao-matrices}
\end{align}
이고 MO coefficient는
\begin{equation}
  C^{\vk\dagger}S^{\vk}C^{\vk}=I
  \label{eq:mo-orthonormality}
\end{equation}
를 만족한다.

\subsection{Restricted density matrix}

Closed-shell occupation $f_{p\vk}\in\{0,2\}$를 사용하면 AO density는
\begin{equation}
  D_{\mu\nu}^{\vk}
  =\sum_p f_{p\vk}C_{\mu p}^{\vk}C_{\nu p}^{\vk *}
  =2\sum_i C_{\mu i}^{\vk}C_{\nu i}^{\vk *}
  \label{eq:restricted-density}
\end{equation}
이다. 이 정의는 occupation factor $2$를 이미 포함한다. 따라서 energy와 response
contraction에서 factor $2$를 다시 넣는지 여부는 사용한 occupied-only tensor가
$D$를 대체하는지에 따라 결정해야 한다.

Matrix trace convention은
\begin{equation}
  \Tr[D^{\vk}X^{\vk}]
  =\sum_{\mu\nu}D_{\nu\mu}^{\vk}X_{\mu\nu}^{\vk}
  \label{eq:trace-convention}
\end{equation}
이다. Complex k-point에서 $D^{\vk}$와 Hermitian operator는 Hermitian이지만 개별
derivative block은 반드시 Hermitian일 필요가 없다. 최종 scalar에는 실수부를 취하되,
중간 단계의 imaginary part를 임의로 버리지 않는다.

\subsection{Index와 derivative notation}

\begin{center}
\begin{tabular}{cl}
\toprule
Index & 의미 \\
\midrule
$\mu,\nu,\lambda,\sigma$ & primary AO \\
$P,Q$ & auxiliary AO \\
$i,j$ & occupied MO \\
$a,b$ & virtual MO \\
$p,q$ & arbitrary MO \\
$A,B$ & atom \\
$x,y$ & Cartesian component \\
$\vk,\vk'$ & sampled k-point \\
\bottomrule
\end{tabular}
\end{center}

핵 좌표 $R_{Ax}$에 대한 total derivative는 superscript $Ax$ 또는
$\partial_{Ax}$로 쓴다. 예를 들어 $h^{Ax}=\partial h/\partial R_{Ax}$다. AO basis가
핵과 함께 움직이므로 matrix derivative와 orbital coefficient response를 구분해야 한다.
Cell은 고정되어 있어 $\mathbf A^{Ax}=0$이고 reciprocal vectors와 discrete mesh도
고정된다.

\subsection{Canonical unit}

\begin{center}
\begin{tabular}{lll}
\toprule
Quantity & Canonical unit & 비고 \\
\midrule
Energy & Hartree (Ha) & public energy \\
Position/lattice & Bohr & PySCF internal \\
Nuclear gradient & Ha/Bohr & $\dd E/\dd R$ \\
Force & Ha/Bohr & $-\dd E/\dd R$ \\
k-point & Bohr$^{-1}$ & Cartesian reciprocal coordinate \\
Plane-wave cutoff & Ha & $G_{\max}^2/2$ \\
\bottomrule
\end{tabular}
\end{center}

입력 XYZ가 Angstrom이어도 cell을 build한 뒤 내부 좌표는 Bohr다. Gradient array를
eV/Angstrom으로 내보낼 때는 energy와 length conversion을 동시에 적용하며, 부호를
바꾸는 것은 unit conversion이 아니라 force 변환이다.
